% The Prose Barrier --- arXiv Preprint
% ACL-format short paper (4 pages + references)
% arXiv: cs.CL / cs.AI

\documentclass[11pt]{article}
\usepackage[preprint]{acl}

\usepackage{times}
\usepackage{latexsym}
\usepackage[T1]{fontenc}
\usepackage[utf8]{inputenc}
\usepackage{microtype}
\usepackage{inconsolata}

\usepackage{booktabs}
\usepackage{amsmath}

\title{The Prose Barrier: Structural Limits of Agent Self-Verification \\
  and a Five-Layer Configuration Integrity Architecture}

\author{Yuhao Lin \\
  Fujian Agriculture and Forestry University \\
  \texttt{3235706088@stu.fafu.edu.cn} \\\\
  \textbf{Independent undergraduate work. No research supervisor. No funding.}}

\begin{document}
\maketitle

\begin{abstract}
AI coding agents are configured through natural-language rules---system prompts,
markdown files, hook registries. Over long conversations, these rules degrade.
The root cause, I find, is structural: an LLM's generation and self-evaluation
share the same decoder $P(\text{token} \mid \text{context}; \theta)$, so the
model cannot independently verify rule compliance---it can only \textit{generate
a claim} that it did. I call this the Prose Barrier.

I present a five-layer architecture organized by position relative to the
Barrier: L0 psychological safety (uncertainty admission), L1 mechanical gates
(filesystem checks bypassing NL entirely), L2 neural gates (logprob probes
detecting constraint penetration without content interpretation), L3 causal
encoding (syllogistic format altering internal representation depth), and L4
drift prediction (trend-based early warning). Deployed over 34 agent sessions
with 11 experiments across 3 architectures (DeepSeek V4 Pro, Qwen3-8B, GLM-4-9B).

Three patterns are visible. First, rules improve compliance (direction
consistent across three architectures; e.g., $+0.38$ above baseline for DeepSeek),
though behavioral scoring reliability was poor
($\kappa=-0.14$). Second, syllogistic format deepens constraint internalization
at the logprob level ($d_z{=}+0.578$, $BF_{10}{=}2.8{\times}10^5$, objective
API-read DV) without changing behavioral output ($\text{IMP}{\approx}\text{SYL}$
across architectures). Third, format effects follow a non-monotonic pattern
($C_1 > C_0 > C_3 > C_2$) governed by output constraint severity, not a slope.

All experiments by a single undergraduate student on one laptop (Dell G15,
RTX 3060). No supervisor. No funding. Code and data at
\url{https://github.com/YuhaoLin2005/hermes-workspace}.
\end{abstract}

\section{Introduction}

A typical Claude Code agent loads 5--10 markdown files defining its behavioral
rules: how to review code, when to run tests, what constitutes a deliverable.
These rules sit in system prompts and hook registries. The agent reads them at
session start. Then it works---sometimes for hours, across dozens of tasks,
generating thousands of tokens.

Whether the agent actually follows its rules is an empirical question.

I ran 34 sessions with a rule-configured agent. Retrospective coding of these
sessions showed rule violations in 55.9\% of sessions (19 of 34) before
mechanical enforcement was wired in. After adding filesystem-level checks
(mtime comparison, regex matching, exit code verification) and running a
150-task controlled experiment across 6 sessions, the violation rate dropped to
0.7\% (1 of 150 tasks). These behavioral numbers come from a scoring protocol
later found to have poor inter-rater reliability ($\kappa=-0.14$; see
Limitations), so they should be treated as directional indicators, not precise
estimates.

The scoring protocol's failure was itself informative: the same structural
constraint that prevents the agent from self-verifying (Section~\ref{sec:barrier})
also prevented the single author from blind-scoring reliably---a mirror of the
Prose Barrier at the meta-level, where independent verification requires a second
person rather than a second prompt.

The core question is \textit{why}. A language model that performs at professional-exam level
can fail to reliably verify that it wrote a file to disk. The answer, I came to
believe, is about architecture, not capability.

The core problem is that when an LLM agent evaluates itself---``Did I follow
rule R on this task?''---its answer comes from the same decoder that may have
already failed to follow R. There is no independent channel. The act of
self-checking and the act of generating share the same token distribution
$P(\text{token} \mid \text{context}; \theta)$. This is what I call the
\textbf{Prose Barrier}: a structural constraint that makes unassisted agent
self-verification unreliable, not because the model is bad at it, but because
the architecture makes genuine independent verification impossible.

I advance three claims:

\begin{enumerate}
\item \textbf{The Prose Barrier is a structural constraint, not an engineering
  bug.} It follows from the autoregressive nature of LLMs and means that
  verification channels must be architecturally independent of generation
  channels. This reframes agent reliability from ``write better prompts'' to
  ``design architectures that work around the Barrier.''

\item \textbf{A five-layer architecture systematically addresses this
  constraint.} The layers are organized by their position relative to the
  Prose Barrier---pre-processing (L0), operating outside it (L1), probing inside
  it without interpreting content (L2), altering internal routing through it
  (L3), and observing it from outside across time (L4).

\item \textbf{Format effects on internal representations are measurable at the
  logprob level, but behavioral claims require independent validation.}
  Logprob-level measurement (L2, $d_z=+0.578$, $BF_{10}=282$k) provides objective
  evidence that syllogistic rule format alters internal constraint representations
  compared to imperative format---without relying on human scoring. Constraint
  gradient experiments reveal a non-monotonic relationship between output
  constraint severity and format effect strength. The behavioral compliance
  numbers (55.9\% $\to$ 0.7\%) are exploratory rather than confirmatory because
  all behavioral experiments were scored by a single unblinded rater; a partial
  blind check yielded $\kappa = -0.14$---below chance. This limitation is
  discussed explicitly, and Section~\ref{sec:whatweneed} specifies what would be
  needed for confirmatory claims.
\end{enumerate}

These claims connect through a single mechanism: a rule that an agent follows
is, at the representation level, a constraint that the decoder respects during
generation. If format changes how deeply a constraint penetrates internal
representations (Claim~3), then format engineering is a lever for rule compliance.
The logprob experiments below measure this penetration at the token-probability
level, bypassing the scoring reliability problem that affects the behavioral
data. The behavioral experiments test whether penetration translates to
observable compliance. The gap between them---format changes representations
but not behavior---is itself the central empirical finding.

\section{The Prose Barrier}
\label{sec:barrier}

\subsection{Definition}

Consider an autoregressive language model $M$ with parameters $\theta$. At
decoding step $n$, the model samples:
\[
t_n \sim P(\cdot \mid t_{1}, \ldots, t_{n-1}; \theta)
\]
Every token the model produces---whether it is ``thinking about what to do,''
``evaluating what it just did,'' or ``taking action''---is drawn from this same
distribution. There is no separate ``auditor'' module with independent access to
ground truth. When the agent is asked ``did you follow rule $R$,'' its answer
is sampled from the same decoder whose state may already encode a failure to
follow $R$.

\textbf{Consequence.} The Barrier cannot be circumvented by prompting. Telling the agent
``be more careful about self-checking'' generates more self-checking
tokens---but those tokens are still drawn from the same distribution. The model
cannot step outside its own generation to independently audit it. This is not a
claim about how well LLMs self-verify; it is a claim about what self-verification
\textit{means} in an autoregressive architecture.

\subsection{The Five-Layer Response}

If the Barrier cannot be removed, it can be worked around. This configuration
integrity architecture addresses agent rule compliance from five positions
relative to the Barrier:

\begin{description}
\item[L0 --- Psychological Safety.] Pre-processes the generation by reframing
  uncertainty. Rather than ``always be helpful and confident,'' the agent is
  told: ``Accuracy $>$ completeness. `I don't know' is a valid output. You will
  not be punished for admitting uncertainty.'' This does not bypass the
  Barrier---it changes what the model generates before the Barrier becomes
  relevant. Tested on 40 binary-choice probes: accuracy preserved ($+0.01$),
  non-ceiling-probe uncertainty admission improved ($r=+0.949$, $n=5$ probes
  below ceiling, 95\% CI $[0.57, 0.996]$).
  \textbf{Status:} preliminary signal only; $n{=}5$ non-ceiling probes,
  a single outlier could drive the correlation.

\item[L1 --- Mechanical Gate.] Operates entirely outside the Barrier. Instead of
  asking the agent ``did you read back the file?'' (NL, goes through the
  Barrier), a shell script compares \texttt{mtime} timestamps (filesystem, no NL
  involved). Violations are detected by regex matching, exit code checking, and
  hook wiring verification---none of which go through the model.
  Deployed over 34 sessions: violation rate dropped from 55.9\% to 0.7\%
  (exploratory; see Limitations for $\kappa=-0.14$ on behavioral scoring).
  19 automated tests (e.g., hook-file mtime verification, SHA256 chain integrity,
  stale output detection via regex, exit-code gate enforcement) operate at the
  filesystem level, bypassing the model entirely; all 19 pass.
  \textbf{Status:} deployed and operational; single-machine only.

\item[L2 --- Neural Gate.] Probes inside the Barrier but avoids content
  interpretation. Rather than reading the agent's output and judging whether it
  ``seems compliant'' (an NL task, unreliable per the Barrier), I measure token
  logprobs at specific decision points. If a syllogistic rule format produces
  higher probability for rule-consistent choice tokens compared to imperative
  format, I have detected format penetration---without interpreting a single
  word of the agent's output. The DV is read directly from the API JSON response
  (\texttt{logprobs=True}, \texttt{top\_logprobs=20}); zero human scoring.
  \textbf{Status:} 40-probe validation complete ($d_z=+0.578$, $BF_{10}=282$k); single
  model (DeepSeek V4 Pro) due to API logprob availability.

\item[L3 --- Causal Encoding.] Alters internal processing through the Barrier.
  The hypothesis: reformulating the same rule from imperative (``verify your
  output'') to syllogistic form (``Write operations can deviate from intent.
  Only a read-back confirms correctness. Therefore, verify after writing.'')
  changes attention routing because the causal chain structure provides a
  template that the model's attention mechanism can follow. Evidence from L2
  logprob measurement, constraint gradient experiments, and convergent
  validation across two independent manipulations.
  \textbf{Status:} multiple experiments completed; causal mechanism relies on
  Pender (2026) for independent mechanistic evidence.

\item[L4 --- Drift Prediction.] Observes L1--L3 state across time, entirely
  outside the Barrier. 12 mechanical features (SHA256 chain integrity, hook
  wiring consistency, stale file detection, etc.) are tracked across sessions.
  When feature patterns deviate from calibration baselines, the system flags
  emerging drift before behavioral failures cascade.
  \textbf{Status:} deployed with 12 features; predictive validation pending.
\end{description}

\begin{table}[t]
\centering
\begin{tabular}{lccc}
\toprule
\textbf{Layer} & \textbf{Position} & \textbf{Mechanism} & \textbf{Evidence} \\
\midrule
L0 Safety  & Pre-barrier & Reframe uncertainty & Accuracy preserved ($+0.01$), \\
            &             & admission             & $r{=}+0.949$ ($n{=}5$ non-ceiling) \\
L1 Gate     & Outside     & Filesystem checks     & 55.9\%$\to$0.7\% violation, 19/19 tests \\
L2 Gate     & Inside      & Logprob probes        & $d_z{=}+0.578$, $BF{=}2.8{\times}10^5$ ($n{=}40$) \\
L3 Encoding & Through     & Syllogistic format    & $d_z$ 0.09--0.60 (non-monotonic) \\
L4 Predict  & Outside     & Drift features        & 12 features deployed, validation pending \\
\bottomrule
\end{tabular}
\caption{Five-layer architecture: position relative to the Prose Barrier,
mechanism, and current empirical support. L1 and L4 operate outside the
Barrier (mechanical checks); L2 probes inside without interpreting content;
L3 alters internal processing through the Barrier; L0 reframes before
generation begins.}
\label{tab:architecture}
\end{table}

\section{Key Experiments}

\subsection{Logprob V3: Does Format Change Internal Representation?}
\label{sec:v3}

An initial attempt ($n=8$ probes) produced $d_z = -0.148$---a near-zero result
in the wrong direction. I wrote a blog post about the null finding~\footnote{DEV.to:
``My Experiment Showed Zero Effect. A Statistician Told Me My Measurement Was
Broken.''}. A reader identified the issue: 4 of 8 probes had comparison tokens
absent from DeepSeek's top-20 logprobs, so my script substituted a $-10.0$
sentinel value---injecting noise that masked the real effect.

I built a mechanical pre-screening pipeline (\texttt{probe\_validator.py}) that
queries the API for top-20 logprobs at the decision token position and verifies
both option tokens are present. Probes that fail are rejected before the main
experiment runs. This is the equivalent of calibrating an instrument before
taking readings.

\textbf{Design.} Within-probe, 3-condition (NO RULES / imperative / syllogistic).
40 probes (10 per category: action, epistemic, structural, meta), all passing
pre-screening. Each probe is a binary choice (A or B) embedded in a scenario
drawn from real Claude Code agent usage. For instance: \textit{``The agent
completed a multi-file refactoring. Should it (A) run the test suite before
declaring completion, or (B) trust the changes and proceed?''} ---a probe for
the action category testing whether the agent prioritizes verification over
speed. The DV is $\log P(A) - \log P(B)$ at
the first token position where both A and B logprobs are available, read directly
from the DeepSeek V4 Pro API response ($T=0.2$, chosen as a compromise:
greedy $T=0$ returns unavailable-logprob sentinels on this API, while $T=1.0$
introduces excessive variance for a repeated-measures design; $T=0.2$ provides
usable logprobs with modest sampling noise).

\textbf{Results.} $n=40$ probes. Mean IMP (imperative-format) differential:
$-2.95$ ($SD=14.67$). Mean SYL (syllogistic-format) differential: $+4.48$
($SD=13.45$). Paired difference
$\bar{x}_{\Delta} = +7.43$, $SE=2.03$, $t(39)=3.65$, $p=0.0008$ (two-tailed).
Cohen's $d_z = +0.578$ (bootstrap 95\% CI on raw logprob differential: $[+3.39, +11.17]$). Bayes factor
$BF_{10} = 282{,}399$ (JZS prior, $r=0.707$), extreme evidence for the
alternative. 32/40 probes (80\%) show direction favoring syllogistic format.

\begin{table}[t]
\centering
\begin{tabular}{lrrr}
\toprule
\textbf{Category} & \textbf{IMP} & \textbf{SYL} & $\mathbf{d_z}$ \\
\midrule
Action      & $-3.66$ & $+1.85$ & $+0.50$ \\
Epistemic   & $-4.23$ & $+3.92$ & $+0.76$ \\
Structural  & $-2.37$ & $+5.56$ & $+0.67$ \\
Meta        & $-1.54$ & $+3.21$ & $+0.38$ \\
\bottomrule
\end{tabular}
\caption{Logprob V3 per-category results. All categories favor syllogistic
format. Category $\times$ Format ANOVA: $F(3,36)=0.26$, $\eta^2=0.02$, n.s.}
\label{tab:v3}
\end{table}

The $SE = 2.03$ yields a 95\% confidence interval of approximately
$[+3.39, +11.17]$ on the raw logprob differential---the effect is robust but
the magnitude is uncertain. Category $\times$ Format interaction was not
significant ($F(3,36)=0.26$, $\eta^2=0.02$). The small $\eta^2$ indicates that the observed interaction effect is near zero
($\eta^2=0.02$); however, with $n=10$ per category, power to detect a medium
interaction ($\eta^2=0.06$) at $\alpha=0.05$ is approximately $0.35$, so the
absence of a detectable interaction should not be interpreted as evidence that
no interaction exists---only that the observed effect size is small.

\subsection{Where Format Helps: L1-Visibility Analysis}
\label{sec:l1vis}

A reviewer of an earlier draft predicted that format effects should concentrate
on rules where mechanical gates cannot reach---format as the ``last line of
defense'' where L1 is absent. I tested the opposite
prediction: format amplifies where L1 provides structural anchors.

\textbf{Classification.} Each of the 40 probes was classified as L1-visible
(violation produces a deterministic, machine-detectable signal) or L1-invisible
(violation concerns judgment quality, with no mechanical ground truth). The
classification uses a three-test criterion: SIGNAL (does violation produce a
deterministic mechanical signal?), ACTION (can a hook verify receipt-of-action?),
CERTAINTY (zero false positives on mechanical signal?). All three must pass for
L1-visible. 22/40 probes (55\%) classified as L1-visible; 18/40 (45\%) as
L1-invisible.

\textbf{Results.} L1-visible probes: $\bar{x}=+9.47$, $s=13.36$, within-group
$d_z=+0.71$, 95\% CI $[+0.35, +1.08]$. L1-invisible probes: $\bar{x}=+5.20$,
$s=12.91$, within-group $d_z=+0.40$, 95\% CI $[+0.02, +0.78]$. The between-group
difference $\Delta d=+0.31$ does not reach significance ($t(38)=1.03$, $p=0.31$)
at $n=40$ for a between-group comparison. A sensitivity analysis reclassifying three ambiguous probes at the
boundary did not change the direction of the effect.

The prediction was wrong. Format helps \textit{more} where L1 already operates
(within-group $d_z=+0.71$ for L1-visible) than where it does not (within-group
$d_z=+0.40$ for L1-invisible), though the between-group difference is not
significant ($t(38)=1.03$, $p=0.31$). The direction is consistent with synergy
(format amplifying existing structural anchors) rather than substitution (format
compensating for missing mechanical enforcement), but $n=40$ is underpowered for
this between-group comparison and the classification was performed by a single
unblinded rater.

\textbf{Limitation.} The classification was done by the same person who ran the
experiment, without a second classifier or inter-rater reliability check. The
three-test criterion was developed to be as mechanical as possible (binary
yes/no for each test), but the boundary between ``receipt-of-action'' and
``receipt-of-diligence'' has gray areas.

\subsection{Constraint Gradient: A Non-Monotonic Pattern}
\label{sec:gradient}

The L2 neural gate measures format effect at the logprob level. But to measure
it, I had to constrain the agent's output---otherwise it produces paragraphs of
natural language where the first-token logprob DV cannot be cleanly extracted.
The constraint text was: ``Output only the letter A or B'' (Chinese original:
``只输出字母A或B'').

This creates a paradox. I am measuring whether syllogistic format produces
deeper reasoning about the constraint, but I am also telling the model to
suppress its reasoning output. The measurement instrument is suppressing the
mechanism it is measuring---what I came to call the \textbf{Prose Barrier of
Measurement}.

This is not a paradox limited to the present design: any behavioral output
constraint on an autoregressive model (forced-choice, short-form, structured
format) risks suppressing the internal processing it aims to measure, because
the same decoder generates both the reasoning and the response.

\textbf{Design.} To test whether this suppression is monotonic (stricter
constraint $\to$ weaker format effect), I ran the 12 P1 probes through 4 levels
of output constraint (labeled C0--C3 to distinguish from architectural layers
L0--L4), 2 formats each, 96 API calls total (DeepSeek V4 Pro, $T=0.2$):

\begin{itemize}
\item \textbf{C0}: No constraint. V3 baseline condition.
\item \textbf{C1}: ``Output only the letter A or B'' (light: standardize output format).
\item \textbf{C2}: ``Output only one letter, no other text, no explanation'' (moderate:
  explicitly suppress explanation).
\item \textbf{C3}: ``Do not output any characters except A or B. No explanation. No punctuation.
  No spaces. No newlines.'' (heavy: maximal output suppression).
\end{itemize}

\textbf{Results.} The monotonic prediction is wrong.

\begin{table}[t]
\centering
\begin{tabular}{lcccc}
\toprule
\textbf{Level} & $\mathbf{n}$ & $\bar{x}$ & $\mathbf{s}$ & $\mathbf{d_z}$ \\
\midrule
C0 (none)     & 11 & $+4.91$  & 15.58 & $\mathbf{0.315}$ \\
C1 (light)    & 12 & $+7.10$  & 11.92 & $\mathbf{0.596}$ \\
C2 (moderate) & 12 & $+1.13$  & 12.40 & $\mathbf{0.091}$ \\
C3 (heavy)    & 12 & $+2.21$  & 7.42  & $\mathbf{0.297}$ \\
\bottomrule
\end{tabular}
\caption{Constraint gradient: format effect rises, crashes, then partially
recovers. C0 $n=11$ due to one probe (degradation-chain FATAL IMP)
having no A/B token at the first position. The C0 $d_z=0.315$ vs.\ the V3
$d_z=0.578$ (Section \ref{sec:v3}) differ because the constraint gradient uses
only the 12 P1 probes for within-experiment comparability, while V3 uses the
full 40-probe pool. Constraint levels C0--C3 are distinct from architectural
layers L0--L4.}
\label{tab:gradient}
\end{table}

The pattern (Table~\ref{tab:gradient}) suggests three qualitatively different
processing regimes rather than a single monotonic function. \textbf{Important
caveat:} no formal statistical test of non-monotonicity was conducted (e.g.,
quadratic contrast or piecewise regression across the four levels). The
three-regime description below is a qualitative interpretation of the observed
means, not a confirmed statistical model, and was developed on the same data
used to evaluate it:

\begin{description}
\item[Optimization ($C_0 \to C_1$):] $d_z$ \textit{rises} from 0.315 to 0.596.
  A light constraint that standardizes output format without telling the model
  to skip reasoning actually \textit{improves} the measurement---clean parseable
  output, still-processed reasoning. This is the best measurement condition for
  logprob-based format effect detection.

\item[Suppression ($C_1 \to C_2$):] $d_z$ \textit{crashes} from 0.596 to 0.091.
  ``No explanation'' directly instructs the model to skip the
  reasoning channel that syllogistic format uses. The measurement constraint
  destroys the mechanism it is measuring. C2 is the worst measurement
  condition---it kills the mechanism while still requiring constrained output.

\item[Rebound ($C_2 \to C_3$):] $d_z$ \textit{recovers} to 0.297 with the
  tightest variance across all levels ($s=7.42$). Under extreme compression,
  the model's system prompt becomes the only remaining behavioral
  differentiator. Syllogistic format's explicit causal structure may be
  embedding more effectively into compressed decision pathways than imperative
  brevity. But the constraint text is also longer (29 vs.~21 characters), and
  this length confound cannot be cleanly separated from severity.
\end{description}

\textbf{Convergent evidence.} The L1-visibility pattern transitions from synergy
to compensation as constraint severity increases:
$\Delta_{V-I}$: $+10.98 \to +0.36 \to -2.03 \to -3.80$. The same reversal
was observed in a separate experiment where multi-scene cognitive load was the
manipulation ($\Delta=-10.08$ in the P1 experiment). Two independent
manipulations producing the same qualitative pattern is suggestive---not
conclusive, given the small $n$ and post-hoc nature of the comparison,
but warrants formal replication.

\textbf{Per-probe variation.} Two probes (self-audit-complexity and
self-audit-logic) survive all four constraint levels with broadly stable format
effects, while the shortest-chain L1-gatable probe (execution-rule-script)
collapses entirely ($f_x$: $+16 \to +4 \to 0 \to -5$). Probes with
moderate-complexity causal chains---not the simplest and not the most
complex---show the most robust format effects. This is consistent with an
inverted-U relationship between chain complexity and format benefit, though
$n=12$ is far too small to estimate a parametric curve.

\subsection{Behavioral Replication Across Architectures}
\label{sec:crossmodel}

\textbf{Rules improve compliance, with caveats.} I ran the same 12 probes in a
condition with all mechanical enforcement hooks removed (the agent's
rule-configuration system, called GateGuard, was set to OFF---all
filesystem-level checks at L1 and behavioral-test hooks were disabled) with a
NO RULES baseline on three models. Compliance rates (NO RULES $\to$ IMP $\to$ SYL):
DeepSeek V4 Pro (MoE): 0.48 $\to$ 0.86 $\to$ 0.83; Qwen3-8B (Dense):
0.79 $\to$ 1.00 $\to$ 0.98; GLM-4-9B (GLM): 0.83 $\to$ 1.00 $\to$ 1.00.
These numbers are behavioral scores from a protocol with $\kappa=-0.14$ and
should be interpreted as exploratory direction indicators, not precise estimates.

\textbf{Format does not change behavior.} IMP compliance $\approx$ SYL compliance
on all three models ($|\Delta| \leq 0.03$). But on DeepSeek, Logprob V3 shows
$d_z=+0.578$. Format changes internal representations without changing behavioral
output---the same dissociation observed in the original GateGuard-OFF experiment
now replicated on two additional architectures. That said, Qwen and GLM show
near-ceiling compliance (IMP/SYL both at $0.98$--$1.00$), so there is little
room for format to create a behavioral difference even if one exists.

\textbf{Cross-model constraint gradient.} I ran the same 4-level constraint
gradient experiment on Qwen3-8B and GLM-4-9B (behavioral measurement only;
logprobs unavailable via SiliconFlow API). Neither model showed any format
effect at any level (GLM: $d_z=0$ everywhere, with perfect 12/12 A-selection
in all conditions). Two possible explanations: (a) 8B--9B models genuinely
lack the processing depth to benefit from syllogistic structure, or (b) the
behavioral measurement method---which already shows $\text{IMP}\approx\text{SYL}$
on DeepSeek---cannot detect format effects at any model scale when ceiling
effects are present. Logprob-level cross-model data would distinguish these
explanations but was not available.

\section{Discussion}

\subsection{What This Work Is and Is Not}

This work has two contributions. The conceptual contribution is reframing agent
reliability from ``how to make them better'' to ``what makes self-verification
structurally impossible''---identifying the Prose Barrier as the root cause of
a practical engineering problem (agent configuration drift). The empirical
contribution is demonstrating that format effects on internal constraint
representations are measurable at the logprob level with an objective,
API-read dependent variable---a method that bypasses the behavioral scoring
problem ($\kappa=-0.14$, see Limitations) that prevents confirmatory behavioral
claims in the current data.

This is not a completed project by the standards of a supervised lab; all
behavioral results should be interpreted as exploratory (see Limitations for
the full list).

\subsection{Beyond the First Token: Context-Fragility of Format Effects}

The Logprob V3 experiment (\S\ref{sec:v3}) measured format effects at the first
decision token in single-scenario probes---optimal conditions: one binary choice,
minimal output constraints. But real agent behavior unfolds across hundreds of
tokens, through multiple decisions, under shifting cognitive demands. Does the
format effect survive these conditions?

\textbf{Multi-scenario collapse.} I ran a follow-up experiment (P1, $n=12$
probes) in which each probe required three simultaneous decisions on distinct
scenarios within a single prompt, rather than the V3 design of one binary choice
per call. The format effect collapsed: Cohen's $d_z$ dropped from 0.58
(single-scene) to 0.19 (multi-scene), a 67\% reduction. Bootstrap 95\% CI for
the multi-scene $d_z$ crosses zero [$-$0.41, $+$0.90]. Critically, the
per-probe format effects in V3 and P1 are \textit{negatively correlated}
($r=-0.65$, 95\% CI [$-$0.89, $-$0.12], $t(10)=-2.71$). The probes most
responsive to syllogistic format in the clean single-scene condition became the
\textit{least} responsive under multi-scene load---a systematic reversal, not
noise-driven attenuation.

\textbf{The Prose Barrier of Measurement.} Two control experiments disentangled
the causes. The multi-scene P1 design confounded two changes: adding multiple
scenarios, and adding a meta-instruction (``Output only the letter A or B'')
required for parseable logprob extraction. Control A (single scene with
meta-instruction, $n=12$) accounted for $\sim$80\% of the total suppression
($d_z$ 1.59$\to$0.31 on the matched probe subset). The meta-instruction alone---
telling the model to suppress explanation and output only a label---
eliminated most of the format effect, even in the single-scene condition.

This is a structural finding about measurement, not a confound to be eliminated.
The meta-instruction is both \textit{measurement necessity} (without it, the
model outputs free-form reasoning and logprob extraction fails) and
\textit{mechanism suppressor} (syllogistic format works by enabling deeper causal
reasoning; the instruction ``do not explain'' tells the model to skip the
reasoning channel that syllogism was designed to enhance). I call this the
\textbf{Prose Barrier of Measurement}: the measurement instrument (output format
constraint) changes the phenomenon being measured (deep rule processing). This is
a corollary of the Prose Barrier---the same decoder generates both the reasoning
and the constrained output, so constraining one necessarily alters the other.

\textbf{A three-parameter model.} These results, combined with the constraint
gradient (\S\ref{sec:gradient}) and cross-model replication
(\S\ref{sec:crossmodel}), suggest that format effect magnitude is a function of
at least three parameters: causal chain complexity (longer chains show larger V3
effects but collapse more under load), available cognitive depth (constrained by
prompt complexity and output format restrictions), and model capacity (8B--9B
models show no format effect; the effect emerges only on the 236B-effective MoE
architecture). The non-monotonic constraint gradient ($C_1 > C_0 > C_3 > C_2$)
and the P1 collapse converge on the same mechanism: format effects require the
deep reasoning channel to be available. When output constraints, cognitive load,
or model capacity limits suppress that channel, syllogistic format loses its
mechanism.

\textbf{What syllogistic format can and cannot do.} The V3 result ($d_z=+0.578$,
$BF_{10}=282$k) establishes that syllogistic format improves constraint
processing under favorable conditions (single decision, reasoning channel open).
The P1 result establishes that this improvement is context-fragile---it does not
survive the transition to complex, multi-decision prompts where the reasoning
channel is suppressed. The behavioral null result ($\text{IMP}\approx\text{SYL}$,
$\Delta\leq 0.03$) across three architectures then indicates that even when the
logprob-level format effect is robust (V3 conditions), it does not translate into
observable behavioral change within the current experimental paradigm.

These boundaries do not invalidate syllogistic format as a tool. They specify
\textit{when} it works: for isolated, high-depth decisions where the model is
permitted to reason. And they specify when it does not: under cognitive load,
with suppressed output channels, or on low-capacity models. This is not a
failure---it is a measurement of boundary conditions, which is what transforms a
phenomenon from ``surprising observation'' into ``engineering tool with known
operating envelope.''

\subsection{Related Work}

The Prose Barrier concept connects to several established results but also
faces challenges from work demonstrating that LLMs can improve through internal
processes. \citet{wei2022cot} showed that chain-of-thought prompting elicits
better reasoning; \citet{wang2022selfconsistency} demonstrated that majority
voting over multiple samples improves accuracy; and \citet{madaan2023selfrefine}
found that iterative self-feedback can refine outputs. These results appear to
challenge the Prose Barrier: if LLMs can self-improve, why can they not
self-verify? The answer is that these methods all operate by changing the
generation distribution $P(\text{token} \mid \text{context})$---longer reasoning
chains, multiple samples, iterative refinement---rather than by providing an
independent verification channel. They improve outputs by altering the
distribution the decoder samples from rather than providing ground-truth
access. This distinction matters: a model can get better at producing
rule-consistent outputs through internal processes, but it still cannot
independently verify that a specific output satisfies the constraint. External
verification (L1) and internal improvement (CoT, self-consistency) are
complementary, not substitutes. The Prose Barrier does not claim self-improvement
is impossible; it claims that ground-truth-independent verification is
structurally unavailable in a single-decoder architecture.

First, \citet{kambhampati2024} established that LLMs require external
verification modules---L1 is an instantiation of this principle for configuration
management rather than planning, extending the LLM-Modulo framework to the
problem of self-referential rule compliance. \citet{bai2022constitutional}
proposed constitutional AI where models self-critique their outputs; the Prose
Barrier argues this self-critique is structurally constrained by the
shared-decoder architecture, and that external verification channels (L1) are
necessary wherever ground truth is mechanically checkable.

Second, \citet{pendereffect2026} provided independent evidence that declarative
vs.\ procedural format alters attention routing in LLMs; my L2/L3 experiments
measure this phenomenon in an applied engineering context with logprob-level
resolution. The non-monotonic constraint gradient extends this line of work by
showing that the relationship between output constraint severity and format
effect is not simply linear. \citet{rath2026agent} documented drift in
multi-agent systems through inter-agent coordination metrics; this work focuses
on single-agent configuration drift through the lens of a structural constraint
(the Prose Barrier) rather than an emergent coordination phenomenon.

Third, \citet{loopengineering2026} proposed deterministic external verification
loops around probabilistic internal LLM loops---L1 is a configuration-domain
instantiation of this principle. \citet{czerwinski2026} introduced the
action-vs-diligence distinction that grounds the L1-visibility classification;
the empirical finding that L1-synergy (amplification, not substitution) is the pattern
adds quantitative evidence to this taxonomy.

Fourth, this work is shaped by two broader critiques of NLP research practice.
\citet{bender2020climbing} argued that language models, trained on form alone,
cannot in principle acquire meaning---a position that supports this caution
against interpreting logprob shifts as evidence of ``deeper understanding''
rather than distributional preference changes. \citet{lipton2018troubling}
identified patterns of overclaiming in ML scholarship, including failure to
distinguish explanation from speculation and the use of suggestive terminology
that carries unearned connotations; I have attempted to avoid these patterns
through explicit hedging, pre-registration disclosure, and quantitative effect
size reporting, though I acknowledge that terms like ``Prose Barrier'' and
``Neural Gate'' are novel coinages that carry rhetorical weight beyond their
empirical grounding.

\subsection{Practical Implications}

Three implications follow from these findings for agent builders. These
implications arise from a pattern that emerged across the full experimental
sequence: the most reliable effects (L1 compliance gains, logprob format effects)
are those measured through architecturally independent channels, while the least
reliable (behavioral compliance scores) required a single unblinded rater. The
reliability hierarchy mirrors the layers' distance from the Prose Barrier.

First,
rules that can be mechanically verified---where a script can detect a violation
of rule $R$ without NL interpretation---are fundamentally more enforceable than
rules that require judgment. Second, the $+0.38$ compliance gain from having
any rules (exploratory estimate) dwarfs format engineering effects in magnitude,
suggesting that investment in L1 mechanical gates yields the highest marginal
return. Third, if measuring format effects via logprobs, light output
constraints preserve the reasoning channel being measured; heavy suppression
destroys it.

\subsection{What Would Be Needed for Confirmatory Claims}
\label{sec:whatweneed}

This paper reports exploratory findings from a single undergraduate student
working without supervision, funding, or multi-model API access. Each of these
requirements for confirmatory evidence is individually achievable: cross-model
logprob replication costs an estimated \$2--5 in API calls, and the probe
validation pipeline (\texttt{probe\_validator.py}) is open-source at
\url{https://github.com/YuhaoLin2005/hermes-workspace}. The barrier to
confirmatory evidence is not resource but design---specifically, the need for
independent blind scoring that the author cannot provide alone.

The following
would be the minimum requirements to upgrade any of the behavioral claims to
confirmatory status:

\begin{enumerate}
\item \textbf{Independent blind scoring.} At least two raters, blinded to
  condition, scoring all behavioral transcripts. Cohen's $\kappa > 0.7$ is the
  conventional threshold for acceptable inter-rater reliability. The current
  $\kappa = -0.14$ (below chance) means the current scoring protocol cannot
  produce confirmatory evidence even in principle.
\item \textbf{Cross-model logprob replication.} Logprob V3 replicated on a
  non-MoE architecture (Claude, GPT-4) with logprobs API support, to rule out
  the possibility that the $d_z=+0.578$ format effect is specific to DeepSeek's
  MoE architecture. Estimated cost: \$2--5 for 40-probe replication.
\item \textbf{Pre-registered constraint gradient hypothesis.} The three-regime
  model (optimization$\to$suppression$\to$rebound) should be pre-registered as a
  planned contrast with a formal test of non-monotonicity (e.g., quadratic
  contrast in repeated-measures ANOVA) on a new holdout probe set distinct from
  the $n=12$ probes used for discovery.
\item \textbf{Second L1-visibility classifier.} The classification of probes as
  L1-visible vs.\ L1-invisible should be performed independently by a second
  rater, with classification agreement reported.
\end{enumerate}

Until these conditions are met, all behavioral findings in this paper should be
treated as hypothesis-generating rather than hypothesis-testing.

\section{Limitations}

The following limitations must be weighed when evaluating this work:

\begin{enumerate}
\item \textbf{Single-rater, unblinded behavioral scoring (fatal).} All behavioral
  experiments were scored by the author with knowledge of experimental
  conditions. A partial blind check on $n=8$ transcripts produced Cohen's
  $\kappa = -0.14$ (95\% bootstrap CI approximately $[-0.8, +0.4]$ with
  $n=8$), indicating the scoring protocol does not meet the threshold for
  confirmatory inference. This low $\kappa$ does not necessarily mean the
  raters disagreed on most cases---with near-ceiling compliance rates in the
  GateGuard-ON condition, extreme base rates can produce low kappa even when
  observed agreement is high (the ``kappa paradox'';~\cite{lipton2018troubling},
  who warn against reporting metrics without context). Regardless of its cause,
  the $\kappa = -0.14$ means that all behavioral compliance numbers
  reported in this paper (including the 55.9\% $\to$ 0.7\% headline result)
  should be interpreted as exploratory, not confirmatory. Independent blind
  scoring by at least two raters achieving $\kappa > 0.7$ is the most critical
  precondition for any claim the current behavioral results support.

\item \textbf{Single-model logprob evidence.} Logprob V3 ($n=40$) and the
  constraint gradient experiment ($n=12 \times 4$ levels) were conducted on
  DeepSeek V4 Pro only. Cross-model family logprob replication (Claude, GPT-4)
  was prevented by API key and logprob availability constraints. The format
  effect findings may be specific to DeepSeek's architecture (MoE with
  CSA/HCA).

\item \textbf{Small sample sizes in key analyses.} The constraint gradient
  uses $n=12$ probes per level; confidence intervals are wide and the
  non-monotonic pattern could shift with additional probes. The three-regime
  model was discovered and described on the same data. L0 Safety Prompt
  correlation ($r=+0.949$) is based on $n=5$ non-ceiling probes; a single
  outlier could drive this correlation. Cross-model behavioral replication
  uses $n=12$ probes with $BF_{01} \approx 2.7$---data favor the null but
  $n=12$ is underpowered to distinguish a genuine null from an imprecise
  measurement.

\item \textbf{Probe pre-screening selection bias.} The \texttt{probe\_validator.py}
  pipeline rejects probes whose option tokens do not appear in DeepSeek's
  top-20 logprobs. This selects probes in which format effects are already large
  enough to push option tokens into the measurable range, potentially inflating
  the pooled format effect estimate ($d_z=+0.578$). The direction and magnitude
  of this bias cannot be estimated from the current data.

\item \textbf{Constraint-length confound.} C2 and C3 constraint texts differ in
  length (21 vs.\ 29 Chinese characters). Severity cannot be fully isolated
  from prompt length in the constraint gradient experiment.

\item \textbf{Post-hoc L1-visibility classification.} The L1-visibility
  classification ($\S$\ref{sec:l1vis}) was performed by the author without
  a second classifier. While the three-test criterion (SIGNAL/ACTION/CERTAINTY)
  was designed to be mechanical, classification ambiguity at boundaries could
  affect the $d_z=+0.71$ vs.\ $+0.40$ comparison (within-group effects; the
  between-group difference $\Delta d=+0.31$ was not significant).

\item \textbf{Lack of pre-registration.} No experimental hypotheses or analysis
  plans were pre-registered. All reported analyses should be considered
  exploratory.

\item \textbf{Single author, single device, no supervision.} All experiments
  were conducted by a single undergraduate student on a single laptop (Dell G15,
  RTX 3060 6GB, 16GB RAM). No research supervisor reviewed the experimental
  design, data analysis, or manuscript. No funding supported this work.
\end{enumerate}

\section{Conclusion}

Agent configuration integrity is not primarily a prompt engineering problem.
It is an architectural problem with a structural root cause: the Prose Barrier.
The five-layer architecture---permit (L0), bypass (L1), detect (L2), encode
(L3), predict (L4)---offers one systematic response. Critically, only L1
constitutes an external verification channel (per \citet{kambhampati2024}'s
LLM-Modulo framework); L2 and L3 operate inside the model and measure internal
representations, not independent verification. The empirical evidence has clear
limitations (see Limitations above) but points in a consistent direction: rules
work, format matters for internal processing but not behavioral output, and the
effect of output constraints on format measurement is non-monotonic, forming
three processing regimes rather than a simple slope. The Prose Barrier of
Measurement---the discovery that the instrument used to measure format effects
can suppress the mechanism being measured---has implications beyond this paper:
it suggests that clean binary-output experimental designs for LLM evaluation may
systematically underestimate format effects.

\section*{Acknowledgments}

This paper grew out of a series of DEV.to articles I wrote while trying to
figure out why my agent kept ignoring its own rules. I am grateful to the
DEV.to community---especially Mike Czerwinski, whose receipt-of-action versus
receipt-of-diligence distinction sharpened the L1-visibility analysis, and
whose Paper-A-versus-Paper-B framing clarified whether format optimization
serves as substitute for mechanical gates or as fallback where gates cannot
reach; Dipankar Sarkar, whose prediction that format effects should
concentrate where mechanical gates cannot reach sent me down the
gate-absent analysis, and whose suggestion to score logprob differentials
at decision tokens rather than averaged across full outputs improved the
measurement; Max Quimby, whose mechanizability-boundary question motivated
the five-layer architecture; the anonymous reader who spotted that my
logprob script was feeding $-10.0$ sentinels into a statistical test; and
everyone who took the time to argue with me in the comments. Those
discussions shaped every experiment in this paper, including the ones
that failed.

The remaining errors, of which I am sure there are many, are mine alone.

\bibliographystyle{acl_natbib}
\bibliography{references}

\end{document}
